\documentclass[11pt,a4paper]{article}
% Preâmbulo comum. Não duplicar pacotes nos módulos.
\usepackage[utf8]{inputenc}
\usepackage[T1]{fontenc}
\usepackage[brazil]{babel}
\usepackage{amsmath,amssymb,amsthm,mathtools}
\usepackage{geometry}
\geometry{a4paper,margin=2.3cm}
\usepackage{microtype}
\usepackage{enumitem}
\usepackage{graphicx}
\usepackage{xcolor}
\usepackage{booktabs}
\usepackage{hyperref}
\hypersetup{colorlinks=true,linkcolor=blue!40!black,urlcolor=blue!50!black}

\setcounter{tocdepth}{2}
\numberwithin{equation}{section}

\theoremstyle{plain}
\newtheorem{teorema}{Teorema}[section]
\newtheorem{proposicao}[teorema]{Proposição}
\newtheorem{lema}[teorema]{Lema}
\newtheorem{corolario}[teorema]{Corolário}

\theoremstyle{definition}
\newtheorem{definicao}[teorema]{Definição}
\newtheorem{exemplo}[teorema]{Exemplo}
\newtheorem{exercicio}{Exercício}[section]

\theoremstyle{remark}
\newtheorem{observacao}[teorema]{Observação}
\newtheorem{contraexemplo}[teorema]{Contraexemplo}
\newtheorem{intuicao}[teorema]{Intuição}

\newcommand{\R}{\mathbb{R}}
\newcommand{\N}{\mathbb{N}}
\newcommand{\Z}{\mathbb{Z}}
\newcommand{\Q}{\mathbb{Q}}
\DeclareRobustCommand{\eps}{\varepsilon}
\DeclareRobustCommand{\vet}[1]{\vec{#1}}
\DeclareRobustCommand{\dif}{\mathrm{d}}
\DeclareRobustCommand{\abs}[1]{\left\lvert #1 \right\rvert}
\DeclareRobustCommand{\norm}[1]{\left\lVert #1 \right\rVert}
\DeclareMathOperator{\grad}{grad}
\DeclareMathOperator{\Div}{div}
\DeclareMathOperator{\rot}{rot}

\newcommand{\base}[1]{\par\medskip\noindent\textbf{Base para relembrar.} #1\par\medskip}
\newcommand{\cuidado}[1]{\par\medskip\noindent\textbf{Quando a hipótese falha.} #1\par\medskip}

\newcommand{\figura}[2]{%
  \begin{figure}[htbp]
  \centering
  \IfFileExists{#1}{%
    \includegraphics[width=0.84\textwidth]{#1}%
  }{%
    \fbox{\parbox{0.8\textwidth}{\centering Figura ainda não gerada:\\ \texttt{#1}}}%
  }
  \caption{#2}
  \end{figure}%
}

\newenvironment{esqueleto}{%
  \par\medskip\noindent\textbf{Esqueleto da demonstração.}\par\nopagebreak
  \begin{enumerate}[leftmargin=1.4em]
}{\end{enumerate}\medskip}

\title{Caderno de prova --- $\eps$-$\delta$}
\author{}
\date{}
\begin{document}
\maketitle
\tableofcontents

\section{Enunciados}

Prove cada afirma\c{c}\~ao a partir das defini\c{c}\~oes. Nos dois primeiros itens, exiba $\delta$ em fun\c{c}\~ao de $\eps$.

\textbf{1.} $\displaystyle\lim_{x\to 2}(3x-1)=5$.

\vspace{6.5cm}

\textbf{2.} $\displaystyle\lim_{x\to 3}x^2=9$.

\vspace{6.5cm}

\textbf{3.} Unicidade do limite. Se $f$ est\'a definida numa vizinhan\c{c}a furada de $a$ e $\lim_{x\to a}f(x)=L$ e $\lim_{x\to a}f(x)=M$, ent\~ao $L=M$.

\vspace{6.5cm}

\textbf{4.} Teorema do confronto. Se existe $\delta_0>0$ tal que $0<\abs{x-a}<\delta_0$ implica $f(x)\le g(x)\le h(x)$, e se $\lim_{x\to a}f(x)=\lim_{x\to a}h(x)=L$, ent\~ao $\lim_{x\to a}g(x)=L$.

\vspace{6.5cm}

\textbf{5.} Se $f$ \'e deriv\'avel em $a$, ent\~ao $f$ \'e cont\'inua em $a$. A derivada \'e $f'(a)=\lim_{h\to 0}(f(a+h)-f(a))/h$, com $h\neq 0$.

\vspace{6.5cm}

\textbf{6.} A fun\c{c}\~ao $f(x)=\abs{x}$ n\~ao \'e deriv\'avel em $0$.

\vspace{6.5cm}

\end{document}
